\subsection{From access to use: what the literature leaves open}

\paragraph{Eligibility, acceptance and the ability to transact.}
Research on Venezuelan migration to Colombia separates public recognition of migrants from access in practice. \citet{ibanez2025} exploit an eligibility cutoff for an earlier regularization program, the \emph{Permiso Especial de Permanencia} (PEP), and examine the consequences for income, public services and financial access. Their design addresses the effect of eligibility for that program, not of the later protection permit or an instant-payment service. In a randomized information study among regularized Venezuelan migrants in Colombia, \citet{barboni2025} investigate how perceptions of social acceptance affect willingness to engage with financial providers. Willingness is informative about a barrier before use, but it is not the same as controlling an account or making a payment. For Bogotá, \citet{martinshields2022} show why the availability of digital technology should not be read as engagement with government services. In finance, the parallel distinction is between having the means to connect, having a provider recognize a credential, and completing a transaction without assistance. The explanatory variables in our survey describe some of those conditions; they do not identify which barrier is decisive.

\paragraph{A remittance is more than the channel in which it arrives.}
Studies of remittances start from different points in the transaction. \citet{zapata_diaspora_2022} examine how Colombian diaspora policy has drawn cross-border household resources into financial markets. The country in that account is an origin state engaging its diaspora; our study observes women living in Colombia with varied international migration histories. \citet{cuadros_verified_2020} study remittance receipt and child welfare in Cali, using household-survey data and a matching design. Their outcomes illuminate the use of household resources but do not capture how a transfer was delivered, whether it stayed digital, or who could direct it. \citet{cirolia_remittance_2022} trace interacting intermediaries, technologies and obligations in two Cape Town remittance-service cases serving the Congolese diaspora. That work suggests what a channel variable can miss: a formal domestic leg can coexist with an informally coordinated cross-border leg. It does not establish which chain was used for any Colombian transfer.

These distinctions determine what we compare. A reported receipt channel describes the latest international remittance. Reported retention or use describes what happens to remittance proceeds in general. Recent domestic digital participation can involve funds from other sources. A cash withdrawal can be consistent with access to a wallet if recipients need cash, face uncertainty about fees, or use cash for obligations for which digital payment is not useful. Conversely, someone who receives cash can make domestic digital payments. Neither outcome measures how much money moved or follows the same transfer through a sequence of stages.

\paragraph{Who controls the funds, and when?}
Access to a financial resource does not establish control over it. \citet{kabeer1999} distinguish resources, agency and achievements. An experiment offering U.S.-based Salvadoran migrants different degrees of control over home-country savings found that savings responded to the arrangement \citep{ashraf2015}. This shows why control matters as a design dimension, not which household member should control funds in our setting. Ethnographic research in Kenya shows that people use digital finance to sustain collective financial relationships as well as individual aims \citep{kusimba_it_2018}. \citet{smyth_challenging_2022} likewise challenges the assumption that linking remittances to formal finance has an unambiguous meaning for Indigenous women's agency in Oaxaca. Evidence on remittances and intimate partner violence among married women in Mexico examines a different outcome in a different population \citep{silveriomurillo_remittances_2022}; it cautions against equating transfers with protection, not against receiving remittances. Joint arrangements can be useful, constrained, or both. Our reported decision arrangements, assistance and recourse items can be read alongside the migrant accounts but do not establish verified coercion, empowerment or safety gains.

\paragraph{Interoperability and the last mile.}
\citet{eaves_rao_verified_2025} separate the functions of identity, payments and data exchange and distinguish technical features from adoption and public-interest safeguards. This matters because a rail can work across participating providers and still require a user to solve other problems: maintaining account access, seeing the fee in advance, confirming the payment and finding someone accountable if it fails. Comparative work on Pix, SINPE Móvil and CoDi treats adoption costs, network effects and coordination among providers as relevant to how usage spreads \citep{argente_verified_2025}. It does not rank those systems on migrant inclusion. A consumer-focused policy brief also makes the case for governance and safeguards rather than treating payment functionality as a proxy for welfare \citep{draper_verified_2024}. We use that brief as a design argument, not evidence that any specific protection worked.

The literature thus gives competing reasons for the same reported behavior. Skills and access may facilitate use, but use can also develop skills. Keeping proceeds digital may reflect convenience, low expected costs, a provider rule, or another person's control of the account. Turning to cash or a trusted intermediary can reflect constraints or a practical choice. Our contribution is to put reported disposition and domestic participation side by side, with their proper bases, and use authored migrant and expert accounts to interpret where frictions can arise. The study does not identify the causal effect of any one of these explanations.
